\ \\ 
Пусть K — кольцо главных идеалов.

\textbf{Утверждение:} Любая цепочка возрастающих идеалов стабилизируется, то есть из того, что $I_1 \subset I_2 \subset I_3 \subset . . . \subset I_n \subset . . .$
(где все $I_k$ — идеалы) следует, что найдётся такое $n$, что $I_n= I_{n+1} = ...$
\\
\Proof
Пусть $(a_0) \subset (a_1) \subset(a_2) \subset...$— возрастающая последовательность идеалов в K. Проверим, что $I := \bigcup\limits_{i=0}^{\infty}(a_i)$ — тоже идеал. Действительно, если b, c $\in$ $I$, то
$b$, $c \in (a_k)$ для некоторого $k \in \N$, \newline поэтому $b$ + $c \in (a_k) \subset I$ и $\forall x \in K : bx, c x \in (a_k) \subset I$.
\newline Поскольку $K$ — кольцо главных идеалов, то $I = (y)$ для некоторого $y\in K$. Но тогда
$y  \in (a_m)$ для некоторого $m \in \N$, поэтому последовательность стабилизируется, начиная с номера $m$, то есть $\forall n \geq m : I_n = I_m$.\EndProof
\bigbreak
\textbf{Утверждение: } В $K$ существует разложение в произведение неразложимых (в смысле 1-й части определения факториального кольца).
\\
\Proof Предположим, что это неверно для элемента $a \in K \backslash \{0\}$. Тогда $a$ можно
представить в виде $a = a_1 b_1$, где $a$, $b_1 \not\in K^* \cup \{0\}$, и у $a_1$ тоже нет разложения. Повторяя этот процесс для $a_1$, $a_2$, . . . , получим последовательность идеалов ($a$) $\subset$ ($a_1$) $\subset$ ($a_2$) $\subset$ ... ,
причем все вложения в ней строгие к силу попарной неассоциированности порождающих
элементов. Это противоречит тому, что последовательность должна стабилизироваться.
Значит, разложение есть у каждого ненулевого элемента.\EndProof
